% ---------------------------------------------------------------------------
% Fachmodul Zahlensysteme (Enrico Ebert)
%
% Inhaltlich abgeleitet aus Abschnitt 4.4 der Dokumentation
% (includes/04_Designentscheidungen/04_04.tex).
%
% Die Kurznotizen stehen in \note{...}. Sie sind standardmäßig unsichtbar.
% Zum Anzeigen in presentation.tex vor \begin{document} ergänzen:
%     \setbeameroption{show notes}
% oder für einen zweiten Bildschirm:
%     \setbeameroption{show notes on second screen=right}
% ---------------------------------------------------------------------------

\addsectionframe{Fachmodul: Zahlensysteme}

% ---------------------------------------------------------------------------
\begin{itemframe}{Zahlensysteme konfigurieren}
  \item Fünf Modi: \textit{Integers}, \textit{Negative}, \textit{Fractional}, \textit{IEEE-754} und \textit{Arithmetik}
  \item Eingaben je nach Basis mit Regex-Filtern validiert, Bitlänge explizit begrenzt
  \item Negative Zahlen in fester 8-Bit-Breite (Wertebereich $-128$ bis $127$, abgestimmt auf die Vorlesung)
  \item Berechnungen durchgängig mit \texttt{BigInt} statt \texttt{Number}: keine Rundungsfehler bei großen Zahlen
\note{
  \textbf{Person: Enrico}

  \textbf{Zeit: 1,5 min}

  \begin{itemize}
    \item Deckt Vorlesungsstoff ab: Basisumrechnung bis IEEE\,754
    \item 5 Modi: Integers, Negative, Fractional, IEEE-754, Arithmetik
    \item Eingaben basisgerecht validiert, Bitlänge begrenzt
    \item Negative: 8 Bit, Wertebereich -128 bis 127 (Vorlesung)
    \item Berechnungen mit \texttt{BigInt} statt \texttt{Number} -- keine Rundungsfehler
    \item Überleitung: Ablauf in der App zeigen
  \end{itemize}
}
\end{itemframe}

% ---------------------------------------------------------------------------
\begin{contentsframe}{Schrittausführung}
\centering
\vspace{3mm}
\resizebox{0.95\textwidth}{!}{%
\begin{tikzpicture}[node distance=15mm and 11mm]
  \node[diagramnode, minimum width=3.4cm, font=\normalsize] (select) {Modus + Eingabe\\wählen};
  \node[diagramnode, right=of select, minimum width=3.4cm, font=\normalsize] (validate) {Eingabe\\validieren};
  \node[diagramnode, right=of validate, minimum width=4.2cm, font=\normalsize] (steps) {Schritt-Objekte\\mit \texttt{BigInt}\\erzeugen};
  \node[diagramnode, right=of steps, minimum width=3.4cm, font=\normalsize] (render) {Rechengitter\\anzeigen};
  \node[diagramnode, below=20mm of validate, minimum width=3.4cm, font=\normalsize] (error) {Fehlermeldung\\mit Hinweis};
  \node[diagramnode, below=20mm of render, minimum width=3.4cm, font=\normalsize] (overflow) {Overflow-\\Warnbanner};

  \draw[diagramarrow] (select) -- (validate);
  \draw[diagramarrow] (validate) -- node[above, font=\small, text=black] {gültig} (steps);
  \draw[diagramarrow] (steps) -- (render);
  \draw[diagramarrow] (validate) -- node[left, font=\small, text=black] {ungültig} (error);
  \draw[diagramarrow] (render) -- (overflow);
  \draw[diagramarrow] (error.west) -| node[pos=0.25, below, font=\small, text=black] {korrigiert} (select.south);
\end{tikzpicture}%
}
\note{
 \textbf{Person: Enrico}

 \textbf{Zeit: 1 min}

 \begin{itemize}
   \item Ablauf in der App, nicht das Rechenverfahren
   \item Modus + Eingabe $\rightarrow$ validieren $\rightarrow$ Schritt-Objekte $\rightarrow$ Rechengitter
   \item Ungültige Eingabe = Hinweis, Nutzer korrigiert
   \item Stufen aufklappbar, i18n-Texte bleiben beim Sprachwechsel erhalten
   \item Overflow-Banner warnt, korrigiert nicht -- didaktisch bewusst
   \item Darstellung passt sich Modus an:
     \begin{itemize}
       \item Wandler: klassische Mathematik-Blöcke
       \item Arithmetik: farbcodiertes Rechengitter
       \item Negative: Vorzeichen, Einer-, Zweierkomplement
       \item IEEE-754: Sign / Exponent / Mantisse getrennt
     \end{itemize}
 \end{itemize}
}
\end{contentsframe}

